% =====================================================================
\chapitre{Étude préalable et état de l'art}
% =====================================================================

Avant de concevoir l'application, il fallait comprendre comment les groupes gèrent
aujourd'hui leurs dépenses et ce que proposent les outils existants. Ce chapitre présente
cette analyse, les concepts d'intelligence artificielle mobilisés, l'étude de faisabilité,
puis les besoins fonctionnels et non fonctionnels qui ont guidé nos choix.

\subsection{Analyse de l'existant}

Avant de concevoir l'application, nous avons observé comment les groupes autour de nous
gèrent aujourd'hui leurs dépenses communes. Dans la plupart des cas, la gestion reste
artisanale : chacun note ce qu'il a payé dans son téléphone ou dans la conversation du
groupe, puis une personne fait le calcul à la fin, souvent dans un tableur. Quand une
application est utilisée, elle ne l'est généralement que par un ou deux membres du groupe,
qui saisissent tout pour les autres. Cette façon de faire présente plusieurs limites :

\begin{itemize}
  \item les dépenses oubliées ou saisies deux fois faussent le résultat final ;
  \item la répartition inégale (une personne qui ne boit pas, une chambre plus chère) est
        difficile à suivre sans erreur ;
  \item personne ne sait, pendant le voyage, où il en est réellement ;
  \item les tickets de caisse sont rarement conservés, alors qu'ils contiennent le montant
        exact, la date et l'enseigne ;
  \item le programme du voyage vit à part, dans une conversation où les propositions se
        perdent et où personne ne sait ce que les autres préfèrent.
\end{itemize}

\subsection{Recueil des besoins}

Ces observations, complétées par les échanges avec notre encadrante et par les retours
d'amis qui ont testé les premières versions, nous ont permis d'identifier les besoins
suivants :

\begin{itemize}
  \item partager une activité avec un groupe et voir les soldes de chacun en temps réel ;
  \item ajouter une dépense en quelques secondes, sans remplir un long formulaire ;
  \item savoir précisément qui doit rembourser qui, avec le moins de virements possible ;
  \item préparer le programme d'un voyage à plusieurs, en votant pour les propositions ;
  \item visualiser ce programme sur une carte et l'organiser pour perdre le moins de temps
        possible en déplacement ;
  \item obtenir de l'aide sur l'application sans quitter l'écran en cours.
\end{itemize}

\subsection{Étude des solutions concurrentes}

Nous avons étudié les applications les plus utilisées pour le partage des dépenses et pour
la préparation des voyages (tableau~\ref{tab:concurrents}).

\begin{table}[H]
\centering
\small
\begin{tabularx}{\textwidth}{|L{3.4cm}|Y|}
\hline
\rowcolor{entete}\textbf{Application} & \textbf{Constat} \\ \hline
Tricount~\cite{tricount} & Partage des dépenses simple et gratuit, groupes sans création de compte
obligatoire. La saisie reste entièrement manuelle et l'application ne propose aucune aide
à l'organisation du voyage. \\ \hline
Splitwise~\cite{splitwise} & Référence du partage de dépenses, avec simplification des dettes. Les
fonctions avancées, dont la lecture des tickets, sont réservées à l'abonnement payant ; aucune
saisie en langage naturel. \\ \hline
Settle Up~\cite{settleup} & Partage des dépenses multi-devises, plutôt orienté vers les groupes
réguliers (colocation, couple). Même logique de formulaire à remplir pour chaque dépense. \\ \hline
Wanderlog~\cite{wanderlog} & Planification de voyage sur carte avec liste des lieux et itinéraire.
Le suivi des dépenses y est secondaire et le choix collectif des lieux n'est pas au cœur de
l'outil. \\ \hline
\end{tabularx}
\caption{Panorama des solutions concurrentes étudiées}
\label{tab:concurrents}
\end{table}

Aucune de ces solutions ne réunit, dans un même outil, le partage des dépenses, une saisie
assistée par l'IA et l'organisation collective du voyage. C'est ce positionnement que
vise \equipay : une seule application pour payer, équilibrer les comptes et décider
ensemble de ce que l'on fait.

\subsection{Concepts clés}

Le tableau~\ref{tab:concepts} définit les principaux concepts mobilisés dans ce projet et
précise leur rôle dans l'application.

\begin{table}[H]
\centering
\small
\begin{tabularx}{\textwidth}{|L{3.4cm}|Y|}
\hline
\rowcolor{entete}\textbf{Concept} & \textbf{Définition et rôle dans le projet} \\ \hline
Grand modèle de langage (LLM) & Réseau de neurones de type \textit{Transformer} entraîné sur de vastes
corpus pour prédire le mot suivant, capable de comprendre et de produire du texte. Nous utilisons
Qwen3-4B-Instruct-2507~\cite{qwen3,qwen3_4b} pour l'assistant, la saisie par phrase et la rédaction
du programme de voyage. \\ \hline
Modèle vision-langage (VLM) & Modèle qui reçoit une image et du texte et répond en texte. Qwen3-VL-4B~\cite{qwen3vl}
lit directement la photo d'un ticket, sans étape d'OCR séparée, et renvoie ses champs au format JSON. \\ \hline
Affinage LoRA & Plutôt que de réentraîner tous les poids, LoRA~\cite{lora} ajoute à certaines couches
deux petites matrices de faible rang, seules apprises. Moins de 1 \% des paramètres sont entraînés,
ce qui rend l'affinage d'un modèle de 4 milliards de paramètres possible sur un seul GPU. \\ \hline
Quantification & Représentation des poids sur moins de bits (ici 4 bits au lieu de 16). La mémoire
nécessaire est divisée par trois à quatre, pour une perte de qualité mesurée et limitée. Elle est
réalisée avec MLX~\cite{mlx}, la bibliothèque d'Apple pour les puces Apple Silicon. \\ \hline
MLOps & Ensemble des pratiques qui rendent un modèle traçable et reproductible : données versionnées,
expériences enregistrées, comparaison des versions et mise en service contrôlée. Nous utilisons
MLflow~\cite{mlflow,mlflow_docs} et son registre de modèles. \\ \hline
Sécurité au niveau des lignes (RLS) & Règles écrites dans PostgreSQL~\cite{postgres_rls} qui filtrent
chaque requête selon l'utilisateur connecté. Un membre ne peut lire que les activités dont il fait
partie, même si l'application était détournée. \\ \hline
\end{tabularx}
\caption{Concepts clés mobilisés dans le projet}
\label{tab:concepts}
\end{table}

\subsection{Étude de faisabilité}

Le tableau~\ref{tab:faisabilite} résume les quatre volets de l'étude de faisabilité menée
avant le développement.

\begin{table}[H]
\centering
\small
\begin{tabularx}{\textwidth}{|L{3.0cm}|Y|}
\hline
\rowcolor{entete}\textbf{Volet} & \textbf{Constat} \\ \hline
Technique & Flutter, Supabase et FastAPI sont des technologies matures et bien documentées. Flutter
permet de viser iOS et Android avec un seul code, et Supabase fournit d'emblée l'authentification,
le stockage des fichiers et le temps réel. \\ \hline
Données & Des jeux publics de tickets annotés existent (CORD, SROIE, WildReceipt), mais aucun ne
couvre les devises qui nous intéressent (dirham, euro). Pour l'assistant et la saisie par phrase,
aucun jeu n'existe : il faut les construire nous-mêmes. \\ \hline
Matériel & Un modèle de 4 milliards de paramètres s'affine en LoRA sur une carte de 16 Go (RTX 5080).
En 4 bits, les trois modèles servis tiennent ensemble dans la mémoire d'un Mac de 18 Go. \\ \hline
Coûts & Aucun coût d'API : les modèles sont ouverts et exécutés en local, Supabase est utilisé dans
son offre gratuite, et les cartes et itinéraires reposent sur des services gratuits (Apple Plans,
OSRM). \\ \hline
\end{tabularx}
\caption{Synthèse de l'étude de faisabilité}
\label{tab:faisabilite}
\end{table}

\subsection{Spécification des besoins fonctionnels}

Les besoins fonctionnels retenus sont regroupés par domaine :

\begin{itemize}
  \item \textbf{Comptes et profil} : inscription par e-mail ou avec Google, profil avec photo,
        thème clair ou sombre, suppression du compte.
  \item \textbf{Activités} : créer une activité (voyage, sortie, quotidien, crédit) avec sa
        devise, rejoindre une activité par code avec validation du créateur, gérer les rôles
        (membre ou spectateur), archiver ou supprimer.
  \item \textbf{Dépenses} : ajouter une dépense avec montant, devise, date, payeur et photo du
        reçu ; la répartir à parts égales, en pourcentages, en montants ou en parts.
  \item \textbf{Soldes} : calculer le solde de chacun en temps réel, proposer le minimum de
        remboursements, afficher graphiques et statistiques.
  \item \textbf{Programme} : pour les voyages et les sorties, lister les lieux à visiter avec
        adresse, date et heure, voter pour ou contre, afficher le programme en liste, en
        calendrier et sur une carte.
  \item \textbf{Intelligence artificielle} : scanner un ticket, décrire une dépense en une
        phrase, organiser le trajet du programme, poser une question à l'assistant.
\end{itemize}

\subsection{Spécification des besoins non fonctionnels}

\begin{figure}[H]
\centering
\begin{tikzpicture}[
  num/.style={circle, fill=white, draw=black!10, minimum size=0.9cm, font=\bfseries, drop shadow={opacity=0.2}},
  titre/.style={font=\large\bfseries, anchor=south},
  texte/.style={font=\scriptsize, align=center, text width=3.9cm, anchor=north},
  lien/.style={draw=#1, line width=0.8pt},
  ]
  % Demi-couronne en six secteurs
  \foreach \i/\c in {0/lagon, 1/violet, 2/bleuclair, 3/orange, 4/recevoir, 5/ambre} {
    \pgfmathsetmacro{\a}{180-\i*30-1.5}
    \pgfmathsetmacro{\b}{180-(\i+1)*30+1.5}
    \fill[\c!85] (\a:1.4) arc[start angle=\a, end angle=\b, radius=1.4] -- (\b:2.9)
                 arc[start angle=\b, end angle=\a, radius=2.9] -- cycle;
    \pgfmathsetmacro{\m}{180-\i*30-15}
    \pgfmathtruncatemacro{\n}{\i+1}
    \node[num] (n\n) at (\m:2.9) {0\n};
  }
  % Libellés : trois à gauche, trois à droite
  \foreach \k/\x/\y/\c/\t/\d in {
      1/-6.1/0.15/lagon/Performance/{Réponse de l'IA en quelques secondes : modèles MLX 4 bits chargés au premier usage.},
      2/-6.1/2.35/violet/Fiabilité/{Repli sur PyTorch, temps estimés si OSRM ne répond pas, bandeau hors ligne.},
      3/-6.1/4.55/bleuclair/Sécurité/{Règles RLS sur toutes les tables, jeton JWT vérifié, reçus privés.},
      4/6.1/4.55/orange/Sobriété/{Aucun service payant ; backend limité à environ 5,8 Go de mémoire.},
      5/6.1/2.35/recevoir/Ergonomie/{Material 3, mode sombre ; l'IA pré-remplit, l'utilisateur valide.},
      6/6.1/0.15/ambre/Maintenabilité/{Migrations versionnées, 89 tests automatiques, modèles tracés dans MLflow.}}
  {
    \node[titre, text=\c!80!black] (t\k) at (\x,\y) {\t};
    \node[texte] at (t\k.south) {\d};
    \ifnum\k<4
      \draw[lien=\c] (t\k.east) ++(0.1,-0.05) -- (n\k);
    \else
      \draw[lien=\c] (t\k.west) ++(-0.1,-0.05) -- (n\k);
    \fi
  }
\end{tikzpicture}
\caption{Synthèse des exigences non fonctionnelles retenues pour l'application}
\label{fig:non_fonctionnels}
\end{figure}

La figure~\ref{fig:non_fonctionnels} synthétise les six exigences non fonctionnelles
retenues : la performance (réponse rapide des modèles malgré une exécution locale), la
fiabilité (l'application continue de fonctionner quand un service externe ne répond pas),
la sécurité (chaque utilisateur ne voit que ses propres activités), la sobriété (aucun coût
et une empreinte mémoire maîtrisée), l'ergonomie (prise en main immédiate) et la
maintenabilité (code testé et modèles traçables).

\subsection{Synthèse et choix retenus}

L'étude de l'existant a montré qu'aucun outil ne combine partage des dépenses, saisie
assistée par l'IA et organisation collective du voyage. Sur cette base, nous avons retenu
les choix résumés dans la figure~\ref{fig:choix}.

\begin{figure}[H]
\centering
\begin{tikzpicture}[
  centre/.style={circle, fill=lagon, text=white, minimum size=2.5cm, font=\bfseries\large, drop shadow={opacity=0.2}},
  choix/.style={rectangle, rounded corners=5pt, draw=#1, line width=0.8pt, fill=white,
                text width=4.1cm, align=center, font=\small, inner sep=6pt, minimum height=1.55cm},
  lien/.style={draw=#1, line width=1pt, -{Latex[length=2.2mm]}},
  ]
  \node[centre] (e) at (0,0) {EquiPay};
  \node[choix=lagon]     (a) at (-5.3, 2.15) {\textbf{Mobile}\\ Flutter, Material 3,\\ carte native Apple Plans};
  \node[choix=violet]    (b) at ( 0.0, 3.25) {\textbf{Données}\\ Supabase : PostgreSQL,\\ RLS, stockage, temps réel};
  \node[choix=orange]    (c) at ( 5.3, 2.15) {\textbf{Service d'IA}\\ FastAPI en local,\\ inférence MLX 4 bits};
  \node[choix=bleuclair] (d) at (-5.3,-2.15) {\textbf{Modèles}\\ Qwen3-4B et Qwen3-VL-4B\\ affinés par LoRA};
  \node[choix=recevoir]  (f) at ( 0.0,-3.25) {\textbf{MLOps}\\ MLflow : suivi des essais,\\ registre et « champion »};
  \node[choix=ambre]     (g) at ( 5.3,-2.15) {\textbf{Itinéraires}\\ OSRM, algorithme d'ordre\\ et MapKit};
  \foreach \n/\c in {a/lagon, b/violet, c/orange, d/bleuclair, f/recevoir, g/ambre}
    \draw[lien=\c] (e) -- (\n);
\end{tikzpicture}
\caption{Synthèse des choix technologiques et fonctionnels retenus}
\label{fig:choix}
\end{figure}

Ces choix constituent le socle technique et fonctionnel détaillé dans le chapitre suivant,
consacré à la conception de l'application.

\subsection{Conclusion du chapitre}

Ce chapitre a posé le cadre du projet à travers l'analyse de l'existant, l'étude des
solutions concurrentes et la définition des concepts clés (modèles de langage et
vision-langage, LoRA, quantification, MLOps). L'étude de faisabilité et la spécification
des besoins fonctionnels et non fonctionnels ont guidé les choix retenus, qui servent de
base à la conception présentée dans le chapitre suivant.
